\section{The finite shared saddle}\label{sa:saddle}
For the entropy \eqref{sa:entropy}, the maximizing loads at positive
entries are
\[
 a_i=\pi^{-1}\arctan(w_i/D),\qquad
                      D=\sqrt{w_1w_2+w_2w_3+w_3w_1}.
\]
They follow from $e'(a)=\log\cot(\pi a)$ and the sum constraint;
zeros are interpreted by continuity. Differentiating in logarithmic
coordinates gives the positive graph form
\begin{equation}\label{sa:laplacian}
 d^2\calQ=\sum_{i<j}\frac{w_iw_j}{2\pi D(w_i+w_j)}
                          (d\log w_i-d\log w_j)^2.
\end{equation}
Terms on zero coordinates are omitted; for a one-variable triangle
the Hessian is zero. The ordinary graph connectivity makes the sum
positive definite after fixing one ordinary scale. This is an explicit
finite-dimensional inference from \eqref{sa:stationarity}, not an
assumption about a Gaussian cycle.

Fix one ordinary radius at its reference value and write
$L_f=\rho_f e^{z_f}$ for the other $r=R-1$ faces. On the fixed local
chart define
\begin{equation}\label{sa:Psi}
 \Psi_N(z,\delta)=\frac1N\sum_e\Loga\mathcal B_{N,e}(\delta_e;L(z))
                         -\frac dN\sum_{f\ne u,v}\Loga L_f(z).
\end{equation}
All constants, including $O(\log N)/N$ terms and the even fold offset,
remain here. Only the explicit $N^{2R}$ is outside this phase.
Proposition~\ref{sa:local-log} gives convergence with derivatives to
\[
 \Psi(z,\delta)=\sum_e\{Q_{b_e}(\delta_e;L)+Q_{c_e}(\delta_e;L)\}
                                       -\sum_{f\ne u,v}\Loga L_f.
\]
At $(0,0)$ its $z$ gradient vanishes by \eqref{sa:stationarity}, and
its Hessian is positive definite by \eqref{sa:laplacian}.
The quantitative holomorphic implicit function theorem now supplies
unique nearby $z_N(\delta)$ on one fixed polydisk, converging with
derivatives. For example, precondition the gradient by the inverse
limiting Hessian: its derivative is uniformly close to the identity
on a fixed ball, and its value at zero tends to zero; contraction on
a smaller ball proves the assertion uniformly in $\delta$.
Set
\begin{equation}\label{sa:AH}
 A_N(\delta)=\Psi_N(z_N(\delta),\delta),\qquad
 H_N(\delta)=\partial_z^2\Psi_N(z_N(\delta),\delta).
\end{equation}
Uniqueness preserves separate evenness and reality.

The coefficient in \eqref{sa:coefficient} has the exact Cauchy measure
$(2\pi)^{-r}d\vartheta$ on this torus. Indeed the numerator has total
degree $dR$, equal to the sum of its ordinary targets; integrating
the common phase gives one. Choosing a reference vertex is a
primitive integer change of torus variables, so no lattice index
multiplies that measure.
One may change the radii and phase origins to $\rho e^{z_N(\delta)}$:
the integrand is a Laurent polynomial and crosses no pole in the
annuli. Near the saddle it is exactly
\[
 (-1)^{dR}N^{2R}e^{NA_N(\delta)}
 \exp\left[-\frac N2\vartheta^{\mathsf T}H_N(\delta)\vartheta
                                      +O(N|\vartheta|^3)\right].
\]
The real part of $H_N$ is uniformly positive definite on a smaller
spectral polydisk. The complementary face torus remains exponentially
small by \eqref{sa:global-angular}: its real radius changes only
within the fixed positive compact, and its phase origin changes by
a small amount. Decrease the spectral chart so this fixed loss exceeds
the change in $\Re A_N$.
Rescale $\vartheta=N^{-1/2}v$. A fixed small neighborhood is dominated
by $e^{-c|v|^2}$; outside it the angular bound applies. Taylor expansion
through fourth order, with uniform analytic derivative bounds, gives
the Gaussian determinant. The cubic term is odd; its square and the
quartic term have relative size $O(N^{-1})$. The remaining Taylor
terms and Gaussian tails obey the same bound after a standard
$|v|\le N^{1/10}$ split. This proves \eqref{sa:H1}, including its
coefficient one and its holomorphic relative error.

At zero, mixed $z,\delta$ derivatives vanish by separate evenness.
Eliminating the face variables therefore does not change the
spectral second derivative. Differentiating \eqref{sa:Q-derivative}
at zero gives
\[
 Q_b''(0)=-2\pi\sqrt{bA+P}/A.
\]
This proves \eqref{sa:hessian}. Each box has a nonmonomial triangle,
so its constant is strictly positive.

